
\chapter{Mid-Term Progress}

This part reports the project at the mid-term review point. It is the state of the
work before the mechanism experiment and the augmentation experiment were run.

\section{System architecture and design}

Data preparation reads the frozen manifests and decodes every image once into a
\texttt{uint8} memory-mapped cache. Train and validation are separate files keyed by
a hash of the path list, so changing the split invalidates the cache instead of
silently reusing it.

Training is a single engine, \texttt{scripts/train\_ddp.py}, parameterised by
\texttt{configs/protocol\_v2.json}. It runs on two GPUs through
DistributedDataParallel, applies augmentation on the accelerator, and writes three
artifacts per run: a best checkpoint, a resumable last checkpoint, and a metrics JSON
carrying the split fingerprint, the learning rate actually applied per epoch, and the
full hyperparameter set.

Evaluation is separate from training on purpose. \texttt{evaluate\_locked\_test.py}
scores PlantVillage splits with bootstrap intervals; \texttt{evaluate\_plantdoc.py}
scores the field dataset through the label map and reports out-of-mapping
predictions rather than hiding them.

Everything reported passes through \texttt{results/registry.json}, which
\texttt{scripts/build\_tables.py} renders into tables. The registry check fails the
build if a document quotes a number not registered against an artifact.

\section{Methodology and implementation}

Four architectures are compared under one frozen protocol:

\begin{table}[H]
\centering
\caption{Architectures compared at mid-term}
\label{tab:arch-mid}
\tabfont
\begin{tabularx}{\textwidth}{@{}lrrY@{}}
\toprule
\textbf{Model} & \textbf{Parameters} & \textbf{Seeds} & \textbf{Note} \\
\midrule
SmallCNN & 398,662 & 3 & trained from scratch \\
ResNet18 & 11,196,006 & 3 & 38-class head \\
EfficientNet-B0 & 4,056,226 & 5 & reference model \\
ConvNeXt-Tiny & 27,849,350 & 3 & \\
\bottomrule
\end{tabularx}
\end{table}

Protocol v2 is AdamW at learning rate $1\times10^{-4}$ with weight decay
$1\times10^{-4}$, global batch 32, 12 epochs, early stopping with patience 3 on
validation macro-F1. SmallCNN is the single documented deviation, using learning rate
$1\times10^{-3}$ and global batch 64, taken from the configuration that produced its
earlier recorded result. Augmentation is horizontal flip, rotation up to 12 degrees,
brightness and per-channel gain, applied on the GPU.

The reference model is trained across multiple seeds so variance is reported rather
than a single run. Macro averages are computed only over classes with non-zero
support, because an earlier 96-image check produced a meaningless macro-F1 of 0.026 by
averaging over absent classes.

\section{Prototype development}

Two things ran at mid-term. The research pipeline is complete and reproducible:
\texttt{scripts/run\_all.py --check} executes six gates, covering the unit tests, the
table build, the registry check, a batch-composition check, a command-contract test,
and failure-path tests for the unattended training notebook.

The second is a static front end in \texttt{app/static/} with four pages: a landing
page, single-image scan with a Grad-CAM slider, batch upload with CSV export, and a
full inference view. It has no build step and opens in a browser.

The front end is not wired to a live model. When no backend URL is configured,
\texttt{app/static/js/api.js} fails fast with code \texttt{demo} and each page shows
recorded checkpoint output instead of inventing a prediction. The service that would
serve live inference is specified but not implemented, so the prototype demonstrates
the interface and the recorded results rather than a deployed classifier.

\section{Dataset and experimental setup}

\begin{table}[H]
\centering
\caption{Datasets and splits}
\label{tab:data-mid}
\tabfont
\begin{tabularx}{\textwidth}{@{}lrrY@{}}
\toprule
\textbf{Dataset} & \textbf{Images} & \textbf{Classes} & \textbf{Splits used} \\
\midrule
PlantVillage & 54,305 & 38 & 36,762 train, 6,544 validation, 10,709 test \\
PlantDoc & 2,578 & 28 & 236 test images, 27 classes scored \\
\bottomrule
\end{tabularx}
\end{table}

Train and validation are disjoint by path and every class appears in both. The test
split is disjoint from both and was opened once for the reference evaluation,
reproducing bit-identical results.

Preprocessing is resize to $224\times224$, conversion to RGB, scaling to $[0,1]$ and
ImageNet normalisation, identical for every model and both datasets. Training ran on
two Tesla T4 GPUs with mixed precision; the fourteen-run sweep completed in about
2.6 hours with no run failures. Evaluation ran on CPU. Versions are pinned in
\texttt{requirements.lock}.

Every run records a dataset fingerprint. All fourteen sweep runs share the same
fingerprint, which is how the project proves they saw the same split rather than
assuming it.

\section{Testing and validation}

The test suite covered split integrity and disjointness, the checkpoint contract
including strict state-dict loading and parameter counts, metric sanity, and the
evaluator output contract.

Two tests caught real errors during this work. ResNet18 had been recorded with the
parameter count of the 1000-class ImageNet model rather than the 38-class one, and a
metrics function signature had been assumed incorrectly.

Beyond unit tests, the training engine has a smoke test covering a distributed run,
the checkpoint contract and the resume fingerprint guard. The unattended notebook has
failure-path tests that execute the real cells against stubbed workers: a failed gate
ships its evidence without starting the sweep, a crashed run does not abort the loop,
a hung run is killed by timeout, and finished seeds are skipped on restart.

\section{Quantitative results and performance analysis}

\begin{table}[H]
\centering
\caption{Architecture comparison under protocol v2, validation macro-F1}
\label{tab:results-mid}
\tabfont
\begin{tabularx}{\textwidth}{@{}lrrY@{}}
\toprule
\textbf{Model} & \textbf{Mean macro-F1} & \textbf{SD} & \\
\midrule
ConvNeXt-Tiny & 0.98951 & 0.00028 & 3 seeds \\
EfficientNet-B0 & 0.98946 & 0.00171 & 5 seeds, reference \\
ResNet18 & 0.98460 & 0.00166 & 3 seeds \\
SmallCNN & 0.91599 & 0.00865 & 3 seeds \\
\bottomrule
\end{tabularx}
\end{table}

The reference model and ConvNeXt-Tiny differ by 0.00005, which is thirty-four times
smaller than the reference model's own seed spread. On this benchmark the
architecture does not decide the outcome.

Field transfer on PlantDoc test, averaged over five protocol-v2 seeds, falls from
0.9928 to 0.2085, a drop of 78.4 percentage points. Scoring all fourteen checkpoints
shows the lab ranking is preserved but the lab gaps are not: ConvNeXt-Tiny 0.3008,
EfficientNet-B0 0.2085, ResNet18 0.1610, SmallCNN 0.0438.

\section{Integration and verification}

Modules are verified at their interfaces rather than assumed to fit. A parity check
builds each model through the package and through the engine and compares parameter
counts and state-dict keys, matching for all four.

The notebooks are generated by builder scripts that embed the engine, the worker and
the protocol byte for byte, and the build fails if an embedded file differs from the
repository version. This exists because a hand-edited notebook cell silently diverged
once already.

The registry is integrated with the documents: \texttt{build\_tables.py --verify}
fails when a claim has no registered run, when an artifact directory is missing, or
when a generated table is stale.

\section{Technical challenges and limitations}

The most expensive problem was a frozen training protocol that looked reasonable and
could not learn. Twelve runs across four architectures came back between 0.07 and
0.95 where the reference model was at 0.99. The cause was not the learning rate,
which is what the configuration suggested. Blocked shuffling was permuting
2{,}048-image blocks on a manifest sorted by class, so the median batch contained
three distinct classes instead of twelve, and BatchNorm normalised using statistics
from three classes. The tell was that the only architecture which trained,
ConvNeXt-Tiny, uses LayerNorm and ignores batch composition. Fixing the shuffle mode
restored reference behaviour.

Limitations recorded at mid-term: the field evaluation rests on 236 images from a
single source with few images per class, so per-class field numbers are indicative.
Part of the gap is likely annotation noise, which this work does not estimate.

\section{Milestone achievement and future work}

Completed at mid-term: a frozen split with a proven disjointness check, a restored
and structurally verified reference checkpoint, four architectures on one engine, a
fourteen-run sweep with zero failures, a complete PlantDoc label map with three
caveats recorded, field evaluation with bootstrap intervals, and calibration analysis
for all four models.

Planned and not started at mid-term: the background-augmentation experiment, a second
independent field dataset with expert labels, and out-of-distribution detection as a
replacement for the confidence threshold.

Part~II reports what happened to each of these.
